\documentclass[10pt,a4paper]{amsart}
\usepackage[margin=19mm]{geometry}
\usepackage{amsmath,amssymb,mathtools}
\usepackage[scr=rsfs,cal=euler]{mathalfa}
\usepackage{xfrac}
\usepackage[unicode,hidelinks,bookmarks=false]{hyperref}
\newcommand{\ie}{i.e.\ }
\newcommand{\cf}{cf.\ }
\newcommand{\resp}{respectively\ }
\newcommand{\Subsection}{\subsection}
\providecommand{\nobreakdash}{\nobreak\hbox{-}}
\setlength{\emergencystretch}{2em}
\multlinegap0pt
\usepackage{tikz}
\usepackage{amssymb}
\usepackage[normalem]{ulem}
\newtheorem{restatable}{Restatable}
\newtheorem{lemma}{Lemma}
\newtheorem{claim}{Claim}
\newtheorem{claimproof}{Claimproof}
\newcommand{\Fc}{
}
\newcommand{\cX}{\mathcal{X}}
\newcommand{\eps}{\varepsilon}
\title{The stochastic block model has the overlap graph property for modularity}
\author{Shankar Bhamidi, David Gamarnik, Remco van der Hofstad, Nelly Litvak, Pawe{\l} Pra{\l}at, Fiona Skerman, Yasmin Tousinejad}
\date{Source: arXiv:2605.10911v1 (CC BY 4.0)}
\begin{document}
\maketitle

\noindent\emph{Source and licence.} The stochastic block model has the overlap graph property for modularity. Version arXiv:2605.10911v1 (CC BY 4.0), distributed by arXiv under the Creative Commons Attribution 4.0 International licence, http://creativecommons.org/licenses/by/4.0/, as recorded in the arXiv licence metadata for this exact version and shown on the abstract page https://arxiv.org/abs/2605.10911v1. Full licence notice: \texttt{licenses/arxiv-2605.10911-CC-BY-4.0.txt}.\par\smallskip

\noindent\textit{Editorial scope.} Counted unit: the article's lemma lem:cycle\_decomposition (source0.tex lines 1220--1229). The article's complete original proof of it is delivered with it (source0.tex lines 1235--1377, one \texttt{proof} environment of 143 lines, which the article places away from the statement). No mathematics is written, repaired or re-derived: the statement and the proof are the article's own lines, in the article's own notation, and no retained line is substituted or re-flowed.  Retained uncounted as the source-local context the proof needs: lines 577--584; lines 634--637.  Nothing was omitted from any retained span, so the package carries no omission record.  No retained line contains a citation, so the wrapper carries no bibliography and no bibliography entry.  No retained line contains a figure environment, an image-inclusion command or drawing code, so no image is delivered and the package declares no asset dependency.  Editorial formatting only, in addition: an AMS \texttt{amsart} preamble that restates the article's own declarations and macros used by the retained lines, namely the environments \texttt{restatable}, \texttt{lemma}, \texttt{claim}, \texttt{claimproof} on separate counters, together with the standard packages those lines need.  The retained environments print in this document as Restatable 1, Lemma 1, Claim 1, Claimproof 1 in order of appearance, whereas the article prints the same items with the numbering of its own full document.  The complete native source of the article as distributed in the arXiv e-print for this exact version is bundled unchanged as \texttt{sources/200328/source0.tex}.
\begin{align}
\nonumber &x_{ij} \geq 0, &\mbox{for all $i,j \in [k]$},\\[0.5em]
%
\label{eq.def.X.orig} &\sum_i x_{ij}  =  1 &\mbox{for all $j \in [k]$},  \\
%
\nonumber &\sum_{i\neq j}x_{ij}  = t, \mbox{ and }\\
\nonumber &\sum_{i}x_{ii} \geq \max_{\sigma} \sum_i x_{\sigma(i)i}.
\end{align}

\begin{restatable}
{proposition}{propbalanced}\label{prop.balanced}
Let $\mathcal{X}_t'$ be defined as in~\eqref{eq.def.X.orig} with the additional restriction that $\sum_j x_{ij} =1$ for all $i\in[k]$. Then $g^{\rm bal}(t)=\max_{X \in \mathcal{X}_t'}g(X)$ is {\em strictly} decreasing on $0<t<k-1$.
\end{restatable}

\begin{lemma}[Cycle decomposition of a nonnegative circulation]
\label{lem:cycle_decomposition}
Suppose $B=(b_{ij})_{i,j\in[k]}$ is a circulation and satisfies $b_{ij}\ge 0$ for $i\neq j$ and $b_{ii}=0$.
Then, for some $m \leq |(i,j) : b_{ij}>0|$,
 there exist \emph{simple directed cycles} 
$C_1,\dots,C_m$ in $G$ and weights $w_1,\dots,w_m>0$ such that, for every $i\neq j$,
\[
b_{ij} \;=\; \sum_{\ell=1}^m w_\ell \,\mathbf{1}\{ \vec{ij}\in C_\ell\}.
\]
\end{lemma}

\begin{proof}[Proof of Proposition~\ref{prop.balanced}]
Fix $0<t_1 < t_2 < k-1$, and let $X^*\in \cX'_k(t_2)$ be such that $g(X^*)=g^{\rm bal}(t_2)$. Note that it suffices to construct $\widetilde{X}\in \cX'_k(t_1)$ such that $g(\widetilde{X})>g(X^*)$ since then we would have
\begin{equation}\label{eq.whyXtilde} g^{\rm bal}(t_1) \geq g(\widetilde{X})>g(X^*) = g^{\rm bal}(t_2).\end{equation}
%
%
The intuition for the proof is that we can regard the non-diagonal entries of $X^*$ as a directed graph; and they form a non-negative circulation. We can then apply Lemma~\ref{lem:cycle_decomposition} to get a directed cycle decomposition $\{(C_\ell,w_\ell)\}_\ell$ such that these sum to the off-diagonal entries of $X^*$. For any directed cycle and $\eps \leq w_\ell$, we may modify an $X\in \cX'(t)$ by deleting an $\eps$ weight from any $x_{ij}$ with $\vec{ij}$ in the cycle, adding $\eps$ weight to $x_{ii}$. We call this a $(C, \eps)$-transfer move, and we will see that this preserves the properties in~\eqref{eq.def.X.orig}, and will thus yield an element of $\cX(t')$ for $t'<t$. We will show that the $g$ value has strictly decreased by this operation. Starting from $X^*$, we use the cycle decomposition to perform a sequence of transfer moves which allows us to construct the required $\widetilde{X}$ in~\eqref{eq.whyXtilde}. Let us now give the details.


We begin by defining the $(C,\eps)$-transfer move on matrix $X$. Given a directed cycle $C$ and $\eps>0$ define $X^{(\eps)}=T_{C,\eps}(X)$ by setting 
\begin{align*}
x^{(\varepsilon)}_{ii} &:= x_{ii}+\varepsilon, \qquad \text{ for all }i\in V(C),\qquad \qquad x^{(\varepsilon)}_{ij} := x_{ij}-\varepsilon, \qquad \text{ for all }\vec{ij}\in \vec{E}(C),
\end{align*}
and leaving all other entries unchanged.

\begin{claim}\label{clm.ginc_stillvalid}
    We \emph{claim} that if $X\in \cX'_k(t)$, $C$ is a directed cycle with $V(C)\subseteq [k]$ and $0< \eps \leq \min_{\vec{ij}\in \vec{E}(C)} x_{ij}$, then, for $X^{(\eps)}=T_{C,\eps}(X)$,
    \begin{itemize}
        \item $g(X^{(\eps)})>g(X)$ and
        \item $X^{(\eps)}\in \cX'_t(t')$ where $t'=t-|C|\eps$.
    \end{itemize}
\end{claim}
\begin{claimproof}
We first prove that $g$ is strictly increased by the transfer move, i.e., $g(X^{(\eps)})>g(X)$. Notice that we may rewrite $g(X)$ using the Frobenius norm. Fix a row $i$, then, since $\sum_j x_{ij}=1$,
\[
\sum_{1\le j<j'\le k} (x_{ij}-x_{ij'})^2
=\frac12\sum_{j,j'}(x_{ij}-x_{ij'})^2
= k\sum_{j=1}^k x_{ij}^2 -\Big(\sum_{j=1}^k x_{ij}\Big)^2
= k\sum_{j=1}^k x_{ij}^2 - 1,
\]
so that
\begin{equation}
\label{eq:g_frob_bal}
g(X)
= k\sum_{i,j} x_{ij}^2 -k 
\end{equation}
Hence, showing $g(X^{(\eps)})>g(X)$ is equivalent to showing $\sum_{i,j}(x^{(\eps)}_{ij})^2> \sum_{i,j}x_{ij}^2$.
Note that only $2|C|$ entries change in the transfer move. Therefore,
\begin{align}
\nonumber
% 
\sum_{i,j} \big(x^{(\varepsilon)}_{ij}\big)^2 - \sum_{i,j} x_{ij}^2 %\\
%\nonumber 
&= \sum_{i \in V(C)} \Big((x_{ii}+\varepsilon)^2-x_{ii}^2\Big)
   +\sum_{\vec{ij}\in \vec{E}(C)} \Big((x_{ij}-\varepsilon)^2-x_{ij}^2\Big)\\
&= 2\varepsilon\Big(\sum_{i \in V(C)} x_{ii} - \sum_{\vec{ij}\in \vec{E}(C)} x_{ij}\Big) + 2|C|\varepsilon^2.\label{eq:diff_after_transfer}
\end{align}
To see that~\eqref{eq:diff_after_transfer} is positive, we argue as follows. Let $\pi\in S_k$ be the permutation that maps $j\mapsto i$ for directed edges $\vec{ij}$ on the cycle, and fixes all other vertices.
Then $\sum_j x_{\pi(j),j}=\sum_{\vec{ij}\in \vec{E}(C)}  x_{ij} + \sum_{i\notin V(C)}x_{ii}$.
However, by the alignment constraint of~\eqref{eq.def.X.orig}, which states that $\sum_{i=1}^k x_{ii} \geq \sum_{i=1}^k x_{\sigma(i)i}$ for all permutations $\sigma$, we have that
$\sum_{i=1}^k x_{ii} \geq \sum_{i=1}^k x_{\pi(i)i}$, so that
\[
\sum_{i\in V(C)}^k x_{ii} = \sum_{i=1}^k x_{ii} - \sum_{i\notin V(C)}^k x_{ii}  \;\ge\; \sum_{j=1}^k x_{\pi(j),j} - \sum_{i\notin V(C)}^k x_{ii} = \sum_{\vec{ij}\in \vec{E}(C)}  x_{ij}.
\]
Thus, \eqref{eq:diff_after_transfer} is at least $2|C|\eps^2>0$, so we have $g\big(X^{(\varepsilon)}\big) > g(X)$, as required.
    




To see the second part of the claim, note first that $x^{(\varepsilon)}_{ij}\ge 0$ for all $i,j\in [k]$ by construction.
Also, row sums remain $1$ since, for rows $i$ with $i\in V(C)$, we decrease one entry by $\eps$ and increase the diagonal
entry by the same amount, and other row sums do not change. The same argument applies to column sums, 
and thus $X^{(\varepsilon)}$ is doubly stochastic. Moreover, 
\[
\sum_{i\neq j} x^{(\varepsilon)}_{ij} \;=\; \sum_{i\neq j} x_{ij} - |C|\varepsilon \;=\; t-|C|\varepsilon=t'.
\]
%
Lastly, for any permutation $\sigma$, {\Fc we note that by the definition of the transfer move, for $i$ such that $\sigma(i)\neq i$, it holds that $x_{ii}^{(\eps)}\geq x_{ii}$ and $x^{(\eps)}_{\sigma(i)i}\leq x_{\sigma(i)i}$. Hence we have
\[\sum_i x^{(\eps)}_{ii}-\sum_{i}x^{(\eps)}_{\sigma(i)i} = \sum_{i \; : \; \sigma(i)\neq i} ( x_{ii}^{(\eps)} - x^{(\eps)}_{\sigma(i)i} ).\]
%
%
 }
%
Thus since $\sum_i x_{ii}\ge \sum_i x_{\sigma(i),i}$ held for $X$, it still holds for $X^{(\varepsilon)}$. We conclude that $X^{(\eps)}$ satisfies the constraints~\eqref{eq.def.X.orig} and $X^{(\eps)}\in \cX'(t')$ which completes the proof of the claim.
\end{claimproof}

We now return to constructing $\widetilde{X}$ to establish~\eqref{eq.whyXtilde}. Recall that $X^*\in\cX_k(t_2)$ is such that $g(X^*)=g^{\rm bal}(t_2)$.
Define $B=(b_{ij})$ by $b_{ij}=x^*_{ij}$ for $i\neq j$ and $b_{ii}=0$; and note that $B$ is a circulation on $[k]$. To see this we need to check that the outflow from any vertex $i\in[k]$ equals the total inflow into $i$. Since the $i$th row and column both sum to $1$ in $X^*$,
\[
\sum_{j\; : \; i\neq j} b_{ij}
\;=\;
\sum_{j\; : \; i\neq j} x^*_{ij}
\;=\;
1-x^*_{ii}
\;=\;
\sum_{j\; :i \neq j} x^*_{ji}
\;=\;
\sum_{j\; :i \neq j} b_{ji},
\]
as required, confirming that $B$ is a circulation. 
%
%
Therefore, by Lemma~\ref{lem:cycle_decomposition}, there exist simple directed cycles $C_1,\dots,C_m$ in $G$
and weights $w_1,\dots,w_m>0$ such that for every $i\neq j$,
\[
b_{ij} \;=\; \sum_{\ell=1}^m w_\ell\,\mathbf{1}\{\vec{ij}\in C_\ell\}.
\]
In particular, the total off-diagonal mass decomposes as
\begin{equation}
\label{eq:cycle_capacity_sum}
t_2 \;=\; \sum_{i\neq j} x^*_{ij}
\;=\; \sum_{i\neq j} b_{ij}
\;=\; \sum_{\ell=1}^m w_\ell \, |C_\ell|.
\end{equation}
%
%
%
Hence, we may choose $L$ such that $\sum_{\ell< L} w_\ell \, |C_\ell| < t_2 - t_1 \le \sum_{\ell\le L}  w_\ell \, | C_\ell|$.
Define $\varepsilon_\ell:=w_\ell$ for $\ell<L$, $\varepsilon_\ell:=0$ for $\ell>L$, and
\[
\varepsilon_L \ :=\ \frac{t_2 - t_1 - \sum_{\ell< L} w_\ell \, |C_\ell| }{r_L},
\]
so that $0<\varepsilon_L\le w_L$ and $\sum_{\ell=1}^L  \varepsilon_\ell \, |C_\ell| =t_2 - t_1$.
%
Now define a sequence of matrices by $X^{(0)}:=X^*$ and, for $\ell=1,\dots,L$,
\[
X^{(\ell)}\ :=\ T_{C_\ell,\varepsilon_\ell}\big(X^{(\ell-1)}\big) \quad\mbox{and}\quad \widetilde{X}:=X^{(L)}
\]
By Claim~\ref{clm.ginc_stillvalid}, each transfer move preserves nonnegativity, double stochasticity, the alignment constraint,
and decreases the off-diagonal mass by exactly $|C_\ell|\, \varepsilon_\ell$. Hence for $\widetilde{X}=X^{(L)}$ the off-diagonal mass is
\[
\sum_{i\neq j} x^{(L)}_{ij}
=
\sum_{i\neq j} x^{(0)}_{ij} - \sum_{\ell=1}^L |C_\ell|\, \varepsilon_\ell
=
t_2-(t_2-1)
=
t_1,
\]
and therefore $\widetilde X\in \cX_k'(t_1)$.
%
Moreover, again by Claim~\ref{clm.ginc_stillvalid}, each transfer move strictly increases $g$, so
\[
g(\widetilde X)
=
g\big(X^{(L)}\big)
>
g\big(X^{(0)}\big)
=
g(X^*),
\]
which establishes~\eqref{eq.whyXtilde}, and we are done.
\end{proof}

\end{document}
